\section*{数据与代码可用性}
随附源码包包含英文主稿、对应中文译稿、两份补充材料、参数记录，以及修订后的有限元、矩阵格式和受控离散计算脚本。在源码包根目录运行 \texttt{python3 reproduction/run\_all.py}，将先重新生成力学基准，再执行残差控制检验、相关跟踪扩展、阻尼扫描和研究框架验证组。上述计算需要 Python~3、NumPy 和 SciPy；重绘论文图表使用 Matplotlib。构建说明另列 PDF 编译步骤。

原稿来源为\href{https://github.com/mop74720-commits/IJRR_CHNN_KAM_v10_20261007}{公开 v10 仓库}，采用提交版本 \href{https://github.com/mop74720-commits/IJRR_CHNN_KAM_v10_20261007/tree/82519542dd057e8fbd6d6879ef33217761e8cbbb}{\texttt{82519542dd05}}，完整标识保存在随附来源记录中。该仓库的 S1/S2 档案继续作为原参数组的计算记录，不作为本次新生成的证据；本次修订也不声称复现全部历史结果。所提供的方向一研究框架决定本稿三个子问题的组织方式。源码包中的来源--章节对照表记录其格式、控制器、收敛论证和数值协议如何进入六节正文。此前的迁移表则单独说明哪些 v10 结果保留为一般理论、哪些已重新计算、哪些改为条件性陈述，以及哪些仅属于历史基准。

\section*{声明}
本文报告数学分析和计算机生成的力学数据。修订后的基准未开展人体、动物、硬件或在轨实验。材料参考值和制造商规格按各自性质引用，不据此暗示器件已经标定。本文件为沿用原论文章节与版式的研究修订稿，不构成发表或录用声明。
